\chapter{Visão geral do sistema}

\eyebrow{arquitetura}

\vspace{4pt}
O Apex é uma única aplicação de desktop nativa. Um núcleo em Go é dono da janela, do
banco de dados e da lógica de negócio; um front end em React desenha a interface
dentro do WebView nativo do sistema operacional; e \emph{sidecars} em Python executam
as partes que precisam de PyTorch ou OpenCV: inferência de modelos, sobreposição
em tempo real e conformação.

\section{Arquitetura}

\begin{center}
\begin{tcolorbox}[colback=codebg,colframe=hair,boxrule=0.6pt,arc=3pt,width=\linewidth]
\ttfamily\small
Binário nativo único (Wails)\\
\hspace*{1em}|\\
\hspace*{1em}+-- Front end React \ (UI escura de instrumento, desenhada pelo WebView)\\
\hspace*{1em}|\hspace*{4em}$\updownarrow$ ponte tipada Go$\leftrightarrow$JS\\
\hspace*{1em}+-- Núcleo Go\\
\hspace*{2em}\ \ \ +-- store \ \ (arquivo SQLite único: a "planilha única")\\
\hspace*{2em}\ \ \ +-- images/ \ (pastas de imagens pareadas em disco)\\
\hspace*{2em}\ \ \ +-- Sidecars Python (iniciados conforme necessário):\\
\hspace*{4em}\ \ \ +-- inference.py \ \ (segmentação de gordura, grau, recorte)\\
\hspace*{4em}\ \ \ +-- realtime.py \ \ \ (sobreposição de gordura ao vivo, subtração de fundo)\\
\hspace*{4em}\ \ \ +-- conformation.py (convexidade integral, analítica)\\
\hspace*{4em}\ \ \ +-- kinect\_capture.py (RGB+depth opcional)
\end{tcolorbox}
\end{center}

\subsection*{Três princípios}

\begin{enumerate}[leftmargin=1.4em,itemsep=4pt]
  \item \textbf{O banco de dados é a fonte da verdade para o pareamento.} Uma imagem
        não carrega seu vínculo no nome do arquivo. O fluxo consiste em selecionar ou
        criar a carcaça e, em seguida, capturar ou importar \emph{dentro} do contexto
        dessa carcaça, de modo que a imagem seja gravada com seu \tele{carcass\_id} já
        definido. A reconstrução pela ordem dos arquivos não é permitida pelo schema.
  \item \textbf{As fontes de captura são intercambiáveis.} A captura por webcam é
        executada no front end via a API de câmera do navegador; o Kinect (depth) é
        executado através de um sidecar Python; ambos são fontes de uma imagem pareada.
  \item \textbf{As medições indicam seu status.} A cor codifica o
        estado: \pillok{validated} (verde) para a segmentação de gordura,
        \pillalert{experimental} (âmbar) para os graus, e o vermelho reservado para
        pareamento quebrado, nunca para uma série de dados.
\end{enumerate}

\section{Pilha de tecnologia}

\begin{center}
\begin{tabularx}{\linewidth}{@{}l l X@{}}
\toprule
\textbf{Camada} & \textbf{Tecnologia} & \textbf{Papel} \\
\midrule
Casca de desktop & Wails v2 (Go) & Janela, embutir o front end, ponte Go$\leftrightarrow$JS \\
Núcleo          & Go 1.23+      & Banco de dados, pareamento, orquestração \\
Banco de dados  & SQLite (pure-Go) & \tele{carcass.db} único: a planilha única \\
Front end       & React 19 + Vite + TypeScript & A interface do instrumento \\
Estilização     & Tailwind 4, shadcn/ui & Tema "Espresso" do Open MCT \\
Inferência      & Python + PyTorch & Modelos de gordura / grau / conformação \\
Visão           & OpenCV, scikit-image, rembg & Recorte, convexidade, remoção de fundo \\
\bottomrule
\end{tabularx}
\end{center}

\section{Onde os dados ficam}

Tudo é armazenado localmente, sob o diretório de dados de aplicação do sistema
operacional:

\begin{lstlisting}
<app data>/carcass-integration/
|-- carcass.db          # o banco de dados único: lotes, carcacas,
|                       #   imagens, avaliadores, graus, analises
`-- images/
    `-- batch_<id>/
        `-- carcass_<tag>/
            |-- rgb_<uuid>.jpg        # a imagem pareada
            `-- analysis/             # saidas dos modelos (mascaras,
                                      #   sobreposicoes, mapas de convexidade)
\end{lstlisting}

\begin{note}[title={O banco de dados único é a "planilha única"}]
O requisito original da pesquisa era uma planilha única, preenchida no momento da
coleta, que registrasse o vínculo imagem--animal. Aqui essa planilha é
\tele{carcass.db}: um arquivo, uma fonte da verdade, portátil e auditável.
\end{note}
